% ============================================================
% 第 12 章  电力系统的无功功率平衡和电压调整
% （教材：何仰赞、温增银《电力系统分析》下册 第 12 章）
% 结构对齐 12.1—12.4；例题、习题不录
% ============================================================

\chapter{电力系统的无功功率平衡和电压调整}

电压是衡量电能质量的重要指标。质量合格的电压应在\textbf{供电电压偏移、电压波动和闪变、
电网谐波、三相不对称程度}四个方面满足国家标准——本章只讨论第一项（电压偏移）。
核心逻辑是一条因果链：\textbf{无功功率平衡决定电压水平}，故调压的本质是\textbf{无功的平衡与再分配}。

\section{电力系统的无功功率平衡}

保证用户处电压接近额定值，是系统运行调整的基本任务之一。
系统的运行电压水平取决于\textbf{无功功率平衡}：各种无功电源的出力应能满足负荷和网络损耗
在额定电压下对无功的需求，否则电压就偏离额定值。以下假定频率维持额定不变。

\subsection{无功功率负荷和无功功率损耗}

\paragraph{无功负荷}
异步电动机在电力系统负荷（特别是无功负荷）中占比很大，故系统无功负荷的电压特性
主要由异步电动机决定。其消耗的无功功率为
\begin{equation}
  Q_M=Q_m+Q_\sigma=\frac{U^{2}}{X_m}+I^{2}X_\sigma ,
\end{equation}
$Q_m$ 为励磁功率，与电压平方成正比（电压较高时因饱和，稍有偏离二次曲线）；
$Q_\sigma$ 为漏抗中的无功损耗——若负载功率不变，电压降低时转差增大、定子电流增大，
故 $Q_\sigma$ 反而上升。两部分叠加的结果是：\textbf{电压降低时，无功负荷先略降后回升}，
这个非线性特性是电压稳定的物理根源（第 15 章）。

\paragraph{无功功率损耗}
\begin{itemize}
  \item \textbf{变压器}：励磁无功 $\Delta Q_0=I_0\%S_N/100$（与电压平方近似成正比）
        与绕组漏抗中的无功 $\Delta Q_T\approx U_s\%S^{2}/(100S_N)$（与负载平方成正比）。
        大型变压器满载时，无功损耗可达其额定容量的 $10\%$ 以上，不能忽略。
  \item \textbf{输电线路}：串联电抗消耗无功
        $\Delta Q_L=(P^{2}+Q^{2})X/U^{2}$；并联电纳发出容性无功（按本书约定取负）
        $\Delta Q_B=-BU^{2}$。$220$~kV 及以下线路 $\Delta Q_L>\Delta Q_B$，
        $500$~kV 及以上因电纳大，轻载时充电功率可能超过消耗，线路呈净容性。
\end{itemize}

\subsection{无功功率电源}

除发电机外，还有四类无功补偿装置：

\begin{table}[htbp]
  \centering
  \caption{无功电源的类型与特性}
  \small
  \begin{tabular}{@{}p{0.16\linewidth}p{0.26\linewidth}p{0.20\linewidth}p{0.28\linewidth}@{}}
    \toprule
    类型 & 无功输出 & 调节 & 特点 \\
    \midrule
    发电机 & 受 $P$-$Q$ 容量曲线限制 & 连续 & 最经济，但受有功出力与稳定限制 \\
    同步调相机 & 可发可吸 & 连续 & 过激发感性无功、欠激发容性无功；损耗大、投资高 \\
    静电电容器 & $Q_C=U^{2}/X_C$ & 分组投切（不连续） & 只能发感性无功；$Q\propto U^{2}$，电压低时出力骤降 \\
    静止补偿器 SVC & 可发可吸 & 快速、平滑 & 由饱和电抗器／TCR／TSC 与固定电容组合；响应快 \\
    静止无功发生器 SVG & 可发可吸 & 极快、范围宽 & 又称 STATCOM；谐波小、占地小 \\
    \bottomrule
  \end{tabular}
\end{table}

\textbf{静电电容器}的无功输出与电压平方成正比：
\begin{equation}
  Q_C=\frac{U^{2}}{X_C},
\end{equation}
这既是它的优点（电压低时少发，有一定自调节性），也是缺点（系统最需要无功时它出力最少）。
实际装置分组接入或切除，实现补偿功率的不连续调节。

\textbf{静止补偿器}由静电电容器与电抗器（经晶闸管或饱和电抗器控制）组合而成：
\begin{itemize}
  \item 饱和电抗器型：工作范围内电压少许变化即引起电流大幅变化；与电抗器串联的电容用于
        斜率校正，改变电容大小可调节补偿器伏安特性的斜率；
  \item TCR 型（晶闸管控制电抗器）$+$ 固定电容：连续可调，但产生谐波，需装滤波器；
  \item TSC 型（晶闸管投切电容器）：不产生谐波，分级投切；
  \item 静止无功发生器（SVG／STATCOM）：响应更快、运行范围更宽、谐波电流小。
\end{itemize}
静止补偿器能在电压变化时快速、平滑地调节无功，满足\textbf{动态无功补偿}的需要；
TCR、TSC 型还可分相补偿以适应不平衡负荷。

\subsection{无功功率平衡}

无功平衡的条件是：在额定电压下，
\keyformula{\sum Q_{G}+\sum Q_{C}=\sum Q_{LD}+\Delta Q_{\Sigma}+Q_{res}}
即无功电源出力之和等于无功负荷、网络无功损耗与备用容量之和。
\begin{itemize}
  \item 无功备用容量一般为最大无功负荷的 $7\%\sim8\%$；
  \item 有功备用按最大有功负荷的百分数留，无功补偿装置按额定容量计；
  \item 330~kV 及以上超高压电网中，线路充电功率可满足无功需求的 $90\%$ 以上，
        须配置\textbf{并联电抗器}吸收过剩充电功率，同时在较低电压等级的配电网配置
        并联电容补偿——两者并存是正常的。
\end{itemize}
\textbf{配置原则}：按无功平衡需要增添补偿容量，并按\textbf{无功就地平衡}原则分配——
小容量、分散的补偿用静电电容器；大容量、配置在系统中枢点的补偿宜用同步调相机或静止补偿器。
采用补偿措施时应统筹兼顾，选用既能发又能吸无功的设备；有大量超高压线路时，
还应考虑发电机\textbf{进相运行}（功率因数超前）的可能性。

\subsection{无功功率平衡和电压水平的关系}

无功平衡与电压水平是同一件事的两种说法：
\begin{itemize}
  \item 无功充足时，系统可维持在额定电压附近运行；
  \item 无功不足时，电压被迫降低——因为负荷的无功需求随电压降低而减少，
        在新的（较低的）电压下重新达到平衡；
  \item 因此\textbf{电压水平是系统无功充裕程度的标志}，而不是可以任意“调”出来的量。
\end{itemize}
这一条常被误解：以为调压只是“把电压调上去”。实际上若无功电源不足，
强行提高电压会使无功更加不足，最终导致电压崩溃。

\corepoint{“无功不足则电压偏低”要说到位：不是“电压低所以无功不足”，而是
\textbf{无功不足 $\to$ 负荷在低电压下减少无功吸收 $\to$ 在新的低电压下平衡}。
所以调压的根本手段是\textbf{增加无功电源或减少无功损耗}，改变变比只是把已有无功
在空间上重新分配，不能增加无功总量。}

\section{电压调整的基本概念}

\subsection{允许电压偏移}

用电设备的额定电压与电网额定电压一致，但实际电压总偏离额定值。
允许偏移（第 1 章）：35~kV 及以上 $\pm5\%$，10~kV 及以下 $\pm7\%$；
对电压质量要求高的用户另有更严的规定。

\subsection{中枢点的电压管理}

由于负荷点很多、不可能对每个负荷点都直接调压，运行中选定一些\textbf{电压中枢点}
（枢纽变电所母线、重要发电厂母线），只要控制中枢点电压在允许范围内，
就能保证由其供电的所有负荷点电压合格。

\begin{table}[htbp]
  \centering
  \caption{中枢点的三种调压方式}
  \small
  \begin{tabular}{@{}p{0.16\linewidth}p{0.34\linewidth}p{0.40\linewidth}@{}}
    \toprule
    方式 & 做法 & 适用 \\
    \midrule
    逆调压 & 最大负荷时升到 $1.05U_N$，最小负荷时降到 $U_N$ & 负荷变化大、线路长、电压损耗大的网络 \\
    顺调压 & 最大负荷时不低于 $1.025U_N$，最小负荷时不超过 $1.075U_N$ & 负荷变化小、线路短 \\
    常调压 & 任何负荷下都维持在同一水平（如 $1.02\sim1.05U_N$） & 负荷变化很小 \\
    \bottomrule
  \end{tabular}
\end{table}

中枢点电压的允许变化范围，由“各负荷点的允许电压范围 $+$ 中枢点到各负荷点的电压损耗”
反推得到：中枢点电压下限 $=$ 电压最低负荷点的允许下限 $+$ 该路径的电压损耗；
上限 $=$ 电压最高负荷点的允许上限 $+$ 其路径的电压损耗。

\subsection{电压调整的基本原理}

以“发电机—升压变—线路—降压变—负荷”的典型供电链为例，负荷点电压为
\keyformula{U_2=\frac{U_G-\Delta U_{\text{网}}-\Delta U_T}{k}}
其中 $U_G$ 为发电机端电压，$\Delta U_{\text{网}}$ 为线路电压损耗，$\Delta U_T$ 为变压器
归算到高压侧的电压损耗，$k$ 为降压变压器变比。

由该式立刻看出调压的三条途径：
\begin{enumerate}
  \item \textbf{改变 $U_G$}：提高发电机端电压（发电机调压）；
  \item \textbf{减小 $\Delta U$}：减少线路与变压器的电压损耗——改变无功分布
        （无功补偿）、减小电抗（串联电容补偿）、提高电压等级；
  \item \textbf{改变 $k$}：改变变压器变比（分接头、有载调压变压器）。
\end{enumerate}

\section{电压调整的措施}

\subsection{发电机调压}

通过调节励磁电流改变发电机端电压。这是\textbf{最经济合理}的调压方式，不需另增设备。
\begin{itemize}
  \item 对孤立运行的发电厂，若发电机到最远负荷点的总电压损耗在允许范围内，
        仅靠发电机调压（如逆调压）即可满足要求；
  \item 多级变压供电系统中，最大负荷时总电压损耗可达 $35\%$、最小负荷时 $15\%$，
        变化幅度达 $20\%$，单靠发电机调压（一般只有 $\pm5\%$）不够；
  \item 多机并列运行时，各机端电压须协调，否则无功分配不合理，故发电机调压一般只作
        \textbf{辅助性措施}。
\end{itemize}

\subsection{改变变压器变比调压}

双绕组变压器在\textbf{高压绕组}设若干分接头（主接头对应 $U_N$），三绕组变压器一般在
高压和中压绕组设分接头，低压绕组不设。改变变比调压即“适当选择分接头”。

\paragraph{降压变压器分接头的选择}
已知高压侧电压 $U_1$、归算到高压侧的变压器电压损耗 $\Delta U_T$、低压侧要求的电压 $U_2$，
变比 $k=U_{1t}/U_{2N}$，则
\begin{equation}
  U_2=\frac{U_1-\Delta U_T}{k}=\frac{\left(U_1-\Delta U_T\right)U_{2N}}{U_{1t}},
\end{equation}
故所需高压侧分接头电压
\keyformula{U_{1t}=\frac{\left(U_1-\Delta U_T\right)U_{2N}}{U_2}}
按最大、最小负荷分别算出 $U_{1t.\max}$、$U_{1t.\min}$，取其算术平均值
$U_{1t.av}$，选最接近的标准分接头，然后\textbf{校验}最大、最小负荷下低压母线电压是否
都在允许范围内。普通变压器只能在停电时改变分接头，运行中只能用同一个固定分接头。

\paragraph{升压变压器分接头的选择}
方法与降压变压器基本相同，但功率方向相反，电压损耗与高压侧电压的关系由“相减”变为“相加”：
\begin{equation}
  U_{1t}=\frac{\left(U_1+\Delta U_T\right)U_{2N}}{U_2}.
\end{equation}
选择发电厂升压变压器分接头时，须保证最大、最小负荷下发电机端电压都在允许范围内，
一般对发电机母线采用逆调压方式。

\paragraph{有载调压变压器}
当“最大与最小负荷下的电压变化幅度超过分接头可调范围（如 $\pm5\%$）”，
或“调压要求的变化趋势与实际相反”（如逆调压）时，靠固定分接头无法满足，须装设
\textbf{带负荷调压}的变压器：
\begin{itemize}
  \item 有载调压变压器：本身具有调压绕组，可在带负荷条件下切换分接头，调节范围较大
        （可达 $\pm8\times1.25\%$ 等）；
  \item 带有附加调压器的加压调压变压器：用附加绕组串入主回路，产生纵、横向附加电压。
\end{itemize}

\subsection{利用无功功率补偿调压}

在负荷点附近装设无功补偿装置，减少流过线路的无功，从而减小电压损耗。
设简单辐射网首端电压 $U_1$ 恒定、末端负荷 $S=P+jQ$、线路阻抗 $R+jX$。
补偿前末端电压
\begin{equation}
  U_2=U_1-\frac{PR+QX}{U_1},
\end{equation}
装设补偿容量 $Q_C$ 后（补偿后末端电压 $U_{2c}$）
\begin{equation}
  U_{2c}=U_1-\frac{PR+\left(Q-Q_C\right)X}{U_1},
\end{equation}
两式相减得所需补偿容量
\keyformula{Q_C=\frac{U_1\left(U_{2c}-U_2\right)}{X}}
上式略去了补偿装置本身的无功电压特性（电容器的 $Q_C\propto U^{2}$），用于初步估算；
精确计算时应迭代修正，因为补偿后电压升高会使电容器实际出力增大。
若已知的是“补偿后要求达到的电压”，也可写成
$Q_C=Q-\dfrac{U_1\left(U_1-U_{2c}\right)-PR}{X}$，两式等价。

\paragraph{补偿装置的选择}
\begin{itemize}
  \item 需要\textbf{容性无功}（提高电压）时用电容器或调相机；
  \item 需要既能发又能吸（电压波动大、需动态补偿）时用静止补偿器或调相机；
  \item 补偿点应靠近无功负荷，遵循“就地平衡、分层分区”原则，避免无功长距离输送。
\end{itemize}

\subsection{线路串联电容补偿调压}

在线路中串联电容器，用电容的负电抗抵消部分线路电抗，减小电压损耗：
\begin{equation}
  \Delta U=\frac{PR+Q\left(X-X_C\right)}{U},
\end{equation}
若要求末端电压由 $U_2$ 提高到 $U_{2c}$，可解出所需串联容抗
\begin{equation}
  X_C=X-\frac{U_1\left(U_1-U_{2c}\right)-PR}{Q}.
\end{equation}
定义\textbf{补偿度}
\begin{equation}
  K_C=\frac{X_C}{X},
\end{equation}
一般取 $K_C<1$（避免谐振与过补偿）。串联电容补偿的优点是提高电压质量的同时
\textbf{提高输电能力与静态稳定极限}（因 $P_{\max}=U_1U_2/(X-X_C)$ 增大），
缺点是需保护（短路时电容器过电压）、可能引发次同步谐振。

\begin{table}[htbp]
  \centering
  \caption{四种调压措施的比较}
  \small
  \begin{tabular}{@{}p{0.14\linewidth}p{0.22\linewidth}p{0.24\linewidth}p{0.30\linewidth}@{}}
    \toprule
    措施 & 手段 & 效果 & 适用与限制 \\
    \midrule
    发电机调压 & 改变 $U_G$ & 直接、经济 & 范围有限（$\pm5\%$），多机并列需协调 \\
    改变变比 & 选分接头、有载调压 & 不增加无功，只重分配 & 固定分接头须停电切换；有载调压投资高 \\
    无功补偿 & 并联电容、调相机、SVC & 减少无功传输，降低电压损耗 & 需有可控无功电源；电容器 $Q\propto U^{2}$ \\
    串联电容 & 减小线路电抗 & 同时提高电压质量与稳定极限 & 需保护，可能谐振；补偿度受限 \\
    \bottomrule
  \end{tabular}
\end{table}

\corepoint{四类措施的本质区别：\textbf{发电机调压与无功补偿是“增加或就地提供无功”，
改变变比是“把已有无功在空间上重新分配”，串联电容是“减少无功的传输损耗”}。
所以无功总量不足时，只有前两类（以及新增补偿）能真正解决问题；只调分接头，
电压只能“此升彼降”。}

\section{调压措施的应用}

\subsection{各种调压措施的合理应用}

调压方案的选择要按“经济性 $+$ 有效性”顺序考虑：
\begin{enumerate}
  \item 首先充分利用\textbf{发电机调压}（不花钱）；
  \item 电压损耗不大而负荷变化较大时，用\textbf{改变变比}（尤其是逆调压场合用有载调压变压器）；
  \item 无功不足或无功传输距离长时，用\textbf{无功补偿}（就地平衡）；
  \item 远距离输电、需要同时提高稳定极限时，用\textbf{串联电容补偿}；
  \item 以上措施组合使用，并校验在各种运行方式下（最大／最小负荷、检修方式）
        都能满足电压要求。
\end{enumerate}

\subsection{各种措施调压效果的综合分析}

不同措施对电压的影响在“空间”和“时间”上不同：
\begin{itemize}
  \item \textbf{空间上}：无功补偿影响其装设点及其邻近节点；改变变比只影响变压器二次侧
        及其下游；发电机调压影响全网（但衰减快）。
  \item \textbf{时间上}：发电机调压、SVC、有载调压变压器可连续、快速调节；
        固定分接头与分组投切电容器只能分级、慢速调节。
  \item \textbf{耦合性}：多种措施同时作用时效果不是简单叠加——例如提高某点电压会增加
        该点负荷的无功需求，抵消一部分调压效果；无功补偿与变压器变比同时调整时，
        须联合优化，否则可能出现“补偿容量增加但电压不变”的情形。
\end{itemize}

\corepoint{调压计算的答题顺序：（1） 由允许电压偏移定中枢点（或负荷点）的电压上下限；
（2） 算出各元件在最大、最小负荷下的电压损耗；（3） 判断单靠发电机调压是否够；
（4） 不够则选分接头（用 $U_{1t}$ 公式，取平均、校验两端）；（5） 仍不够则加无功补偿
（用 $Q_C$ 公式）；（6） 最后校验所有运行方式。每一步都要“算出数、再校验”，
不能只写公式。}
